% ============================================================================
%  Rozdział 2. Tło teoretyczne i wykorzystane technologie
% ============================================================================
\chapter{Tło teoretyczne i wykorzystane technologie}
\label{rozdz:podstawy}

Przed rozpoczęciem implementacji konieczne było uporządkowanie wiedzy o~dwóch obszarach.
Pierwszy dotyczy samej gry: jakie cechy decydują o~tym, że rozgrywka zręcznościowa jest
czytelna i~angażująca. Drugi obejmuje narzędzia programistyczne, czyli silnik Unity,
jego model obiektowy i~system fizyki, oraz biblioteki, które pozwalają uporządkować kod
ponad to, co oferuje sam silnik. W~rozdziale omówiono kolejno: pętlę rozgrywki gier
zręcznościowych, model komponentowy i~cykl życia skryptów w~Unity, system fizyki,
wstrzykiwanie zależności z~biblioteką Zenject, programowanie reaktywne z~biblioteką R3,
asynchroniczność opartą na UniTask, wzorzec MVVM oraz wzorce projektowe stosowane
w~grach.

\section{Pętla rozgrywki w grach zręcznościowych}

Jesse Schell opisuje grę jako system, w~którym gracz podejmuje decyzje w~ramach reguł,
a~gra natychmiast odpowiada na te decyzje \cite{schell}. W~grach zręcznościowych podstawową
jednostką tego systemu jest krótka pętla: gracz ocenia sytuację, wykonuje akcję, obserwuje
jej skutek i~na tej podstawie planuje następną. W~opisywanej grze pętla ta trwa od kilku
do kilkunastu sekund: gracz wybiera kierunek i~siłę strzału, kula toczy się, odbija od
ścian i~zbiera monety, po czym gracz koryguje tor kolejnym strzałem.

Scott Rogers zwraca uwagę, że gracz uczy się zasad gry przede wszystkim przez informację
zwrotną, a~nie przez tekst \cite{rogers}. Każde zdarzenie powinno mieć widoczny skutek:
zebrana moneta znika z~animacją, licznik monet pulsuje, kula po zderzeniu ze ścianą
zostaje na chwilę zgnieciona i~zmienia kolor, a~utrata zdrowia jest widoczna na pasku
życia. Z~punktu widzenia architektury oznacza to, że jedno zdarzenie w~logice gry musi
trafić do kilku niezależnych odbiorców. To bezpośrednio uzasadnia zastosowanie wzorca
obserwatora i~programowania reaktywnego, omówionych w~dalszej części rozdziału.

Gatunek cechuje także powtarzalność: gracz wraca do tego samego poziomu, aby poprawić
wynik, więc gra powinna mierzyć postęp i~przechowywać najlepsze rezultaty, a~kolejne
poziomy powinny zwiększać trudność bez zmiany zasad sterowania.

\section{Silnik Unity i jego model komponentowy}

Unity jest silnikiem gier rozwijanym przez Unity Technologies, w~którym logikę aplikacji
pisze się w~języku C\# \cite{unity-manual}. W~projekcie wykorzystano wersję 2022.3 LTS,
czyli wydanie o~przedłużonym okresie wsparcia, zalecane do projektów, w~których
stabilność jest ważniejsza od dostępu do najnowszych funkcji.

\subsection{Obiekty gry, komponenty i prefabrykaty}

Podstawowym elementem sceny jest obiekt gry (\klasa{GameObject}), który sam w~sobie nie ma
żadnego zachowania. Zachowanie i~wygląd nadają mu dołączone komponenty: transformacja,
siatka i~materiał, zderzacz, ciało sztywne lub skrypt użytkownika dziedziczący po klasie
\klasa{MonoBehaviour}. Obiekt skonfigurowany w~edytorze można zapisać jako prefabrykat
(\emph{prefab}), czyli szablon, z~którego w~czasie działania gry tworzy się dowolną liczbę
kopii. Scena jest plikiem zawierającym hierarchię obiektów; gra składa się z~kilku scen
ładowanych kolejno przez menedżera scen.

Model komponentowy sprzyja kompozycji zamiast dziedziczenia, ale ma istotne ograniczenie:
instancje komponentów tworzy silnik, a~nie programista. Klasa dziedzicząca po
\klasa{MonoBehaviour} nie może mieć konstruktora z~parametrami, więc zależności muszą być
ustawiane po utworzeniu obiektu. W~projekcie rozwiązano to przez podział klas na dwie
grupy. Cienkie komponenty \klasa{MonoBehaviour} odbierają zależności z~kontenera i~zdarzenia
silnika, a~właściwą logikę realizują zwykłe klasy C\# (w~kodzie oznaczone przyrostkiem
\kod{Logic}), tworzone przez konstruktor.

\subsection{Cykl życia skryptów}

Silnik wywołuje metody komponentów w~ściśle określonej kolejności \cite{unity-order}.
Najważniejsze z~nich przedstawiono na rysunku~\ref{rys:cykl-zycia}. Metoda \kod{Awake}
jest wywoływana raz, zaraz po utworzeniu obiektu, \kod{OnEnable} przy każdej aktywacji,
a~\kod{Start} przed pierwszą klatką, w~której obiekt jest aktywny. Następnie w~każdej
klatce wykonywane są \kod{Update} i~\kod{LateUpdate}, a~niezależnie od nich, ze stałym
krokiem czasu, \kod{FixedUpdate}, w~którym należy modyfikować ciała sztywne. Dezaktywacja
i~zniszczenie obiektu wywołują \kod{OnDisable} oraz \kod{OnDestroy}; w~tych metodach
projekt zwalnia subskrypcje zdarzeń.

\begin{figure}[!htb]
  \centering
  \begin{tikzpicture}[diag, node distance=4mm and 7mm]
    \node[modul, minimum width=2.4cm] (awake) {Awake};
    \node[modul, minimum width=2.4cm, right=of awake] (enable) {OnEnable};
    \node[modul, minimum width=2.4cm, right=of enable] (start) {Start};
    \node[akcent, minimum width=2.6cm, below=12mm of awake] (fixed) {FixedUpdate\\[-1pt]{\scriptsize krok 0{,}02 s, fizyka}};
    \node[akcent, minimum width=2.6cm, right=of fixed] (update) {Update\\[-1pt]{\scriptsize raz na klatkę}};
    \node[akcent, minimum width=2.6cm, right=of update] (late) {LateUpdate\\[-1pt]{\scriptsize kamera, interfejs}};
    \node[modul, minimum width=2.4cm, below=12mm of update] (disable) {OnDisable};
    \node[modul, minimum width=2.4cm, right=of disable] (destroy) {OnDestroy};
    \draw[strz] (awake) -- (enable);
    \draw[strz] (enable) -- (start);
    \draw[strz] (start.south) -- ++(0,-4mm) -| (fixed.north);
    \draw[strz] (fixed) -- (update);
    \draw[strz] (update) -- (late);
    \draw[strz] (late.east) -- ++(5mm,0) |- ($(fixed.south)+(0,-3mm)$) -- (fixed.south);
    \node[font=\scriptsize, anchor=west] at ($(late.east)+(5mm,-0.55)$) {następna klatka};
    \draw[strzp] (update.south) -- (disable.north);
    \draw[strz] (disable) -- (destroy);
    \begin{scope}[on background layer]
      \node[draw=black!40, dashed, rounded corners, fit=(fixed)(update)(late), inner sep=5mm,
            label={[font=\scriptsize]below left:pętla gry}] {};
    \end{scope}
  \end{tikzpicture}
  \caption{Uproszczona kolejność wywołań metod cyklu życia komponentu \klasa{MonoBehaviour}; Źródło: Opracowanie własne na podstawie \cite{unity-order}}
  \label{rys:cykl-zycia}
\end{figure}

Kolejność ta ma bezpośrednie konsekwencje dla projektu. Zenject wstrzykuje zależności
przed wywołaniem \kod{Awake} i~\kod{Start}, więc w~metodzie oznaczonej atrybutem
\kod{[Inject]} wolno korzystać wyłącznie z~przekazanych argumentów, a~nie z~innych
komponentów, które mogły jeszcze nie zostać zainicjalizowane. Subskrypcje zakładane
w~\kod{Start} muszą z~kolei zostać zwolnione w~\kod{OnDestroy}, ponieważ ponowne
wczytanie sceny niszczy wszystkie jej obiekty.

\subsection{System fizyki}

Do symulacji fizyki trójwymiarowej Unity wykorzystuje silnik NVIDIA PhysX
\cite{unity-physics}. Obiekt podlega symulacji, jeśli ma komponent \klasa{Rigidbody};
jego kształt kolizyjny określa zderzacz (\klasa{SphereCollider}, \klasa{BoxCollider},
\klasa{CapsuleCollider} lub \klasa{MeshCollider}). Zderzacz może działać w~dwóch trybach.
W~trybie zwykłym obiekty nie przenikają się, a~silnik wywołuje metody
\kod{OnCollisionEnter}, \kod{OnCollisionStay} i~\kod{OnCollisionExit}, przekazując punkty
styku oraz wektory normalne powierzchni. W~trybie wyzwalacza (\emph{trigger}) obiekty
przenikają się swobodnie, a~silnik zgłasza jedynie zdarzenia \kod{OnTriggerEnter}
i~\kod{OnTriggerExit}. Monety, apteczki, pola przyspieszające i~meta są w~grze
wyzwalaczami, a~ściany, przeciwnicy i~przeszkody zwykłymi zderzaczami.

Ciało sztywne może być dynamiczne, czyli poruszane przez siły i~kolizje, albo kinematyczne,
czyli przesuwane wyłącznie przez skrypt metodami \kod{MovePosition} i~\kod{MoveRotation}.
Ruchome ściany i~obracające się belki są ciałami kinematycznymi: poruszają się po zadanym
torze, ale przy tym poprawnie wypychają kulę, z~którą się zderzają. Symulacja jest
wykonywana ze stałym krokiem czasu, w~projekcie równym domyślnej wartości 0{,}02~s,
co odpowiada 50~krokom fizyki na sekundę niezależnie od liczby klatek obrazu.

Ważnym narzędziem jest również rzutowanie promieni (\emph{raycasting}), czyli wyznaczanie
pierwszego zderzacza przeciętego przez półprostą \cite{unity-raycast}. Funkcja
\kod{Physics.Raycast} przyjmuje parametr \klasa{QueryTriggerInteraction}, który pozwala
pominąć wyzwalacze. W~grze wykorzystano to do sprawdzania, czy gracz dotknął kuli,
nawet gdy kula znajduje się wewnątrz pola przyspieszającego lub apteczki.

\subsection{Universal Render Pipeline, system wejścia i zestawy kompilacji}

Do renderowania wykorzystano Universal Render Pipeline (URP) w~wersji 14, czyli potok
renderowania przeznaczony dla szerokiego zakresu urządzeń, w~tym telefonów \cite{urp}.
Obsługę wejścia zrealizowano za pomocą pakietu Input System, w~którym akcje gracza
definiuje się niezależnie od urządzeń, a~następnie przypisuje do nich konkretne kontrolki
\cite{inputsystem}. Jedna akcja wciśnięcia może więc reagować zarówno na lewy przycisk
myszy, jak i~na dotknięcie ekranu.

Kod gry podzielono na zestawy kompilacji (\emph{assembly definitions}) \cite{unity-asmdef}.
Każdy plik \kod{.asmdef} definiuje osobną bibliotekę z~jawnie zadeklarowaną listą
zależności. Ma to dwie zalety. Po pierwsze, zmiana w~jednym module powoduje ponowną
kompilację tylko tego modułu i~zależnych od niego, a~nie całego projektu. Po drugie,
kompilator nie pozwala odwołać się do klasy z~zestawu, który nie został zadeklarowany jako
zależność, więc niezamierzone powiązania pomiędzy modułami są wykrywane na etapie
kompilacji, a~nie przeglądu kodu.

\section{Wstrzykiwanie zależności i biblioteka Zenject}
\label{sec:di}

Wstrzykiwanie zależności (\emph{Dependency Injection}) polega na tym, że obiekt nie tworzy
ani nie wyszukuje swoich zależności, lecz otrzymuje je z~zewnątrz, najczęściej przez
konstruktor \cite{seemann}. Miejsce, w~którym aplikacja składa graf obiektów, nazywa się
korzeniem kompozycji (\emph{composition root}). Dzięki temu klasy zależą od abstrakcji
i~nie znają sposobu tworzenia swoich współpracowników, co ułatwia wymianę implementacji
i~testowanie. Alternatywą powszechną w~projektach Unity są singletony oraz lokalizator
usług; Nystrom zwraca uwagę, że oba rozwiązania ukrywają zależności i~utrudniają
określenie, który kod korzysta z~danego obiektu \cite{nystrom}.

Zenject (rozwijany również pod nazwą Extenject) jest kontenerem wstrzykiwania zależności
przeznaczonym dla Unity \cite{zenject}. Jego konfiguracja opiera się na trzech pojęciach.
Kontekst (\klasa{ProjectContext}, \klasa{SceneContext}) określa zasięg kontenera:
kontekst projektu istnieje przez cały czas działania aplikacji, a~kontekst sceny tylko do
jej zamknięcia, przy czym kontener sceny dziedziczy powiązania kontenera projektu.
Instalator (\klasa{MonoInstaller}) to klasa, w~której metodzie \kod{InstallBindings}
rejestruje się powiązania typów z~instancjami. Powiązanie (\emph{binding}) zapisuje się
płynnym interfejsem, np. \kod{Container.Bind<T>().FromInstance(obiekt).AsSingle()}.

\begin{figure}[!htb]
  \centering
  \begin{tikzpicture}[diag, node distance=5mm and 6mm]
    \node[akcent, minimum width=12.4cm, minimum height=9mm] (proj)
      {\textbf{ProjectContext} (prefab w~katalogu Resources, istnieje przez cały czas działania gry)\\
       BootstrapInstaller: \kod{LoadLevel} (nazwa poziomu do wczytania)};
    \node[modul, below left=9mm and -2.9cm of proj, minimum width=3.6cm] (menu)
      {\textbf{SceneContext}\\ \textit{Main Menu}\\[1pt]
       BallMenuInstaller\\ WallMenuInstaller};
    \node[modul, right=of menu, minimum width=3.6cm] (load)
      {\textbf{SceneContext}\\ \textit{Loading}\\[1pt]
       LoadingInstaller\\ \phantom{x}};
    \node[modul, right=of load, minimum width=3.6cm] (lvl)
      {\textbf{SceneContext}\\ \textit{Level 1, 2, 3}\\[1pt]
       Ball-, Coins-, Enemies-,\\ Ground-, UI-, Finish-, \dots};
    \draw[strz] (proj.south -| menu.north) -- (menu.north);
    \draw[strz] (proj.south -| load.north) -- (load.north);
    \draw[strz] (proj.south -| lvl.north) -- (lvl.north);
    \node[font=\scriptsize, below=2mm of load] {kontenery scen dziedziczą powiązania kontenera projektu};
  \end{tikzpicture}
  \caption{Hierarchia kontekstów Zenject w~opracowanej grze; \zrodloWlasne}
  \label{rys:konteksty}
\end{figure}

Zenject obsługuje wstrzykiwanie przez konstruktor, pola, właściwości oraz metody.
W~przypadku komponentów \klasa{MonoBehaviour} stosuje się wstrzykiwanie przez metodę
oznaczoną atrybutem \kod{[Inject]}, która pełni rolę konstruktora. Kontener wstrzykuje
zależności do wszystkich komponentów obecnych w~scenie podczas jej wczytywania, natomiast
obiekty tworzone później muszą powstawać przez metodę \kod{DiContainer.InstantiatePrefab},
aby ich komponenty również otrzymały zależności. Strukturę kontekstów zastosowaną w~grze
przedstawia rysunek~\ref{rys:konteksty}.

\section{Programowanie reaktywne i biblioteka R3}
\label{sec:reaktywne}

Programowanie reaktywne traktuje zdarzenia jako strumienie wartości rozłożone w~czasie.
Podstawą jest wzorzec obserwatora \cite{gof}: obiekt obserwowany (\emph{observable})
publikuje kolejne wartości, a~zarejestrowani obserwatorzy otrzymują je w~kolejności
publikacji. Biblioteki z~rodziny ReactiveX rozszerzają ten wzorzec o~operatory, takie jak
filtrowanie, łączenie i~przekształcanie strumieni, oraz o~jednolity sposób kończenia
subskrypcji \cite{reactivex}.

W~ekosystemie Unity przez wiele lat standardem była biblioteka UniRx \cite{unirx}.
W~lutym 2024 roku jej repozytorium zostało zarchiwizowane, a~opis projektu odsyła do
biblioteki R3 tego samego autora jako rozwiniętej wersji UniRx. R3, którego wersja 1.0
ukazała się 16 lutego 2024 roku, jest implementacją Reactive Extensions napisaną od
podstaw z~myślą o~współczesnej platformie .NET i~silnikach gier \cite{r3}. W~projekcie
wykorzystano trzy jej typy:
\begin{itemize}[leftmargin=1.9cm, itemsep=1pt]
  \item \klasa{Subject<T>}, czyli strumień, do którego kod wysyła wartości metodą
        \kod{OnNext}; w~grze reprezentuje zdarzenia jednorazowe, np. zebranie monety
        (\klasa{Subject<int>} z~wartością monety) lub dotarcie do mety
        (\klasa{Subject<Unit>}, gdzie \klasa{Unit} oznacza brak wartości),
  \item \klasa{ReactiveProperty<T>}, czyli właściwość, która przechowuje bieżącą wartość
        i~powiadamia obserwatorów o~jej zmianie; w~grze przechowuje m.in. tekst licznika
        monet, prędkość kuli i~wynik poziomu,
  \item subskrypcję zwracaną przez metodę \kod{Subscribe} jako obiekt \klasa{IDisposable};
        wywołanie \kod{Dispose} kończy odbieranie zdarzeń i~zwalnia referencje.
\end{itemize}

Istotną różnicą względem zwykłych zdarzeń C\# (\kod{event}, \kod{Action}) jest to, że
subskrypcja jest obiektem, który można przechowywać, grupować i~zwalniać razem z~obiektem
właściciela. Wyrażenie lambda przypisane do zdarzenia C\# trudno później odpiąć, ponieważ
każde jego wystąpienie tworzy nowy delegat; w~przypadku R3 wystarczy zachować zwrócony
obiekt i~wywołać na nim \kod{Dispose}. W~projekcie oba mechanizmy współistnieją: zdarzenia
C\# są używane wewnątrz pojedynczych modułów, a~strumienie R3 tam, gdzie zdarzenie trafia
do kilku niezależnych odbiorców lub przekracza granicę modułu.

\section{Programowanie asynchroniczne z wykorzystaniem UniTask}

Animacje i~opóźnienia w~Unity tradycyjnie realizuje się za pomocą korutyn, czyli metod
zwracających \klasa{IEnumerator}, wznawianych przez silnik w~kolejnych klatkach. Korutyny
wymagają jednak obiektu \klasa{MonoBehaviour}, nie zwracają wyniku i~nie obsługują wyjątków
w~naturalny sposób. Biblioteka UniTask udostępnia własny typ zadania zgodny ze składnią
\kod{async}/\kod{await}, zrealizowany jako struktura, dzięki czemu typowe operacje nie
alokują pamięci na stercie \cite{unitask}. Oczekiwanie jest powiązane z~fazami pętli gry
silnika, więc instrukcja \kod{await UniTask.WaitForFixedUpdate()} wznawia metodę
w~kolejnym kroku fizyki, podobnie jak zrobiłaby to korutyna.

W~grze UniTask wykorzystano do animacji zgniatania kuli, zmniejszania zebranych monet,
płynnej zmiany licznika monet, wypełniania paska ładowania oraz odmierzania czasu animacji
śmierci przeciwnika. Metody niezwracające wyniku mają typ \klasa{UniTaskVoid}
i~są uruchamiane wywołaniem \kod{Forget()}, a~dłuższe animacje przyjmują token anulowania
(\klasa{CancellationToken}), który pozwala przerwać je w~dowolnym momencie.

\section{Wzorzec Model--View--ViewModel}

Wzorzec Model--View--ViewModel (MVVM) został opisany w~2005 roku przez Johna Gossmana
na potrzeby aplikacji tworzonych w~technologii WPF \cite{gossman}. Dzieli on warstwę
prezentacji na trzy role \cite{ms-mvvm}. Model przechowuje dane i~logikę dziedziny,
widok odpowiada wyłącznie za wyświetlanie, a~model widoku przekształca dane modelu do
postaci gotowej do prezentacji i~udostępnia je widokowi przez mechanizm powiązań. Widok
zna model widoku, ale model widoku nie zna widoku, a~model nie zna żadnego z~nich
(rysunek~\ref{rys:mvvm}).

\begin{figure}[!htb]
  \centering
  \begin{tikzpicture}[diag, node distance=6mm and 19mm]
    \node[modul, minimum width=3.0cm, minimum height=1.5cm] (v) {\textbf{View}\\ komponent TextMeshPro,\\ suwak, obraz};
    \node[akcent, minimum width=3.0cm, minimum height=1.5cm, right=of v] (vm) {\textbf{ViewModel}\\ \kod{ReactiveProperty<string>}\\ \kod{Subject<Unit>}};
    \node[modul, minimum width=3.0cm, minimum height=1.5cm, right=of vm] (m) {\textbf{Model}\\ stan gry,\\ np. liczba monet};
    \draw[strz] ([yshift=3mm]v.east) -- node[above, font=\scriptsize]{subskrypcja} ([yshift=3mm]vm.west);
    \draw[strzp] ([yshift=-3mm]vm.west) -- node[below, font=\scriptsize]{nowa wartość} ([yshift=-3mm]v.east);
    \draw[strz] ([yshift=3mm]vm.east) -- node[above, font=\scriptsize]{subskrypcja} ([yshift=3mm]m.west);
    \draw[strzp] ([yshift=-3mm]m.west) -- node[below, font=\scriptsize]{zdarzenie zmiany} ([yshift=-3mm]vm.east);
  \end{tikzpicture}
  \caption{Zależności we~wzorcu MVVM w~wariancie opartym na strumieniach reaktywnych; \zrodloWlasne}
  \label{rys:mvvm}
\end{figure}

Unity nie posiada mechanizmu powiązań danych dla systemu interfejsu uGUI, który został
użyty w~projekcie. Rolę powiązań pełnią więc subskrypcje strumieni R3: widok subskrybuje
właściwości reaktywne modelu widoku i~w~reakcji na zmianę ustawia tekst lub wartość
suwaka. Model widoku może przy tym realizować logikę prezentacji niezależną od komponentów
silnika, np. stopniowe zwiększanie wyświetlanej liczby monet lub zaokrąglanie prędkości do
jednego miejsca po przecinku.

\section{Wzorce projektowe stosowane w grach}

Oprócz wzorców omówionych wyżej w~projekcie wykorzystano kilka rozwiązań typowych dla
programowania gier \cite{nystrom}. Fabryka (\emph{Factory}) skupia w~jednym miejscu
tworzenie obiektów określonego rodzaju; w~grze fabryki tworzą kulę, monety, przeciwników
i~efekty cząsteczkowe, korzystając z~kontenera Zenject. Pula obiektów (\emph{Object Pool})
polega na jednorazowym utworzeniu zestawu obiektów i~ich wielokrotnym użyciu zamiast
tworzenia i~niszczenia w~trakcie rozgrywki. Ogranicza to liczbę alokacji, a~w~środowisku
z~automatycznym odśmiecaniem pamięci także liczbę przerw wywołanych pracą odśmiecacza,
które w~grze objawiają się chwilowym spadkiem płynności.

Trzecim rozwiązaniem są komponenty znacznikowe (\emph{marker components}), czyli puste klasy
dziedziczące po \klasa{MonoBehaviour}, które pozwalają rozpoznać rodzaj obiektu bez
porównywania nazw lub tagów. Sprawdzenie \kod{TryGetComponent(out BallMarker \_)} jest
odporne na literówki i~zmiany nazw obiektów w~scenie. Czwartym jest sterowanie danymi
(\emph{data-driven design}): parametry rozgrywki, takie jak siła tarcia, obrażenia
przeciwników czy siła pól przyspieszających, są przechowywane w~plikach konfiguracyjnych,
a~nie w~kodzie, dzięki czemu strojenie rozgrywki nie wymaga ponownej kompilacji.
